\documentclass{article}
\usepackage{my_style}
\usepackage{amsmath}
\usepackage{nicematrix}


\begin{document}
\mytitle[Expedition im Dschungel]{Musterlösung}

\section*{1. Möglichkeit: Analysis}
Effektiv wollen wir die Distanz zwischen dem besten Startpunkt und dem Tempel minimieren.

Sei $T=
\begin{pNiceMatrix}
    4\\3
\end{pNiceMatrix}
$. 
Auf der Straße befinden wir uns in Abhängigkeit von $x$ bei 
$
\begin{pNiceMatrix}
    x\\\frac 27x
\end{pNiceMatrix}
$

Die Distanz ergibt sich durch die Funktion: $f:\mathbb{R}\to\mathbb{R}$ mit
\begin{align*}
    f(x)=&\sqrt{\left(\frac 27x - 3\right)^2+(x-4)^2}\\
    =&\sqrt{\frac4{49}x^2-\frac{12}7x+9+x^2-8x+16}\\
    =&\sqrt{\frac{53}{49}x^2-\frac{68}7x+25}\\
\end{align*}
Hierdrauf können wir jetzt ganz normal Extremstellen bestimmen:

\begin{align*}
    &0\stackrel!=f'(x)=\underbrace{\frac{1}{2\sqrt{\frac{53}{49}x^2-\frac{68}7x+25}}}_{\not = 0 \forall x}
    \cdot \frac{106}{49}x-\frac{68}7\\
    \Longleftrightarrow &0=\frac{106}{49}x-\frac{68}7\\
    \Longleftrightarrow &\frac{68}7=\frac{106}{49}x\\
    \Longleftrightarrow &x=\frac{68}7\cdot\frac{49}{106}=\frac{238}{53}=4.\overline{4905660377358}\\
\end{align*}

Wir rechnen jetzt nicht mehr die Krümmung aus, wir haben eine nach oben geöffnete Parabel, die hat
nur eine Extremstelle und das ist ein Minimum. \\cjVd:w
Genauer wäre es so: $f$ ist eine Normfunktion und daher 
streng konvex, außerdem ist das Problem unbeschränkt und damit folgt, dass $f$ genau 
ein Minimum hat und dieses auch sicher angenommen wird.

\newpage
\section*{2. Möglichkeit: Lineare Algebra}
Gegeben ist die Funktion \(s(x)=\frac27x\) und der Punkt \(T=(4, 3)\).
Wir schreiben \(s\) in Vektordarstellung als \[S=t\begin{pmatrix}7\\2\\\end{pmatrix}, t\in \mathbb{R}.\]
Jetzt betrachten wir alle Vektoren von \(S\) zu \(T\) mit
\[V=x\begin{pmatrix}7\\2\\\end{pmatrix} - \begin{pmatrix}4\\3\\\end{pmatrix}=
\begin{pmatrix}7t-4\\2t-3\end{pmatrix}.\] In dieser Menge suchen wir jetzt den Senkrechten Vektor
zu \(S\). Dafür berechnen wir zuerst das Skalarprodukt
\[\begin{pmatrix}7t-4\\2t-3\end{pmatrix}\cdot\begin{pmatrix}7\\2\end{pmatrix}=49t-28+4t-6,\] setzen 
es dann mit \(0\) gleich und lösen nach \(t\) auf:
\[0=49t-28+4t-6\Leftrightarrow t=\frac{34}{53}.\]
Wenn wir jetzt dieses \(t\) in unsere Geradengleichung \(S\) einsetzen erhalten

\[\frac{34}{53}\begin{pmatrix}7\\2\\\end{pmatrix}=
\begin{pmatrix}\frac{238}{53}\\\frac{68}{53}\\\end{pmatrix}=
\begin{pmatrix}4.\overline{4905660377358}\\1.\overline{2830188679245}\\\end{pmatrix}
.\]
Die \(x\)-Koordinate von diesem Vektor ist die Lösung: \(4.\overline{4905660377358}\).\\


\newpage
\section{3. Möglichkeit: Lineare Algebra (aber mit Formel)}

Wir können auch den Ortsvektor vom Tempel auf die Richtung der Straße projezieren:

$T=(4, 3)^T, S=(1, \frac27)^T$

\begin{align*}
    p &= \frac{\langle T,S \rangle}{\langle S,S \rangle} S
      = \frac{4+ \frac67}{1+\frac4{49}} \begin{pNiceMatrix}1\\ \frac 27\end{pNiceMatrix}
      = \frac{238}{53} \begin{pNiceMatrix}1\\ \frac 27\end{pNiceMatrix}
      =\begin{pNiceMatrix}\frac{238}{53}\\\frac{68}{53}\\\end{pNiceMatrix}
      =\begin{pNiceMatrix}4.\overline{4905660377358}\\1.\overline{2830188679245}\\\end{pNiceMatrix}
\end{align*}


\end{document}


